% Agent 20: Pages 416--419 (PDF pages 38--39)
% Sections: 5.2 (end) -- vertical pressure balance, sources of viscosity,
%           radiative transport, opacity
% Note: These pages continue and largely complete Section 5.2.

%% ============================================================
%% Page 416, left column (continuation of S5.2)
%% ============================================================

(see \S2.5). Thus, the heat generated per unit area is
\begin{equation}
2h\epsilon = (2h\bar{t}_\phi^{\;r})(-2\varpi\dot{\phi})
= \frac{3\dot{M}_0 M}{4\pi r^2}\left[1-\beta\!\left(\frac{r_1}{r}\right)^{\!1/2}\right];
\tag{5.2.19}
\end{equation}
and the heat generated between radii $r_1$ and $r_2$ is
\begin{align}
\int_{r_1}^{r_2}\! 2h\epsilon\;2\pi r\,dr
&= \frac{3}{2}\,\dot{M}_0\!\left(\frac{M}{r_1}-\frac{M}{r_2}\right)
   \!\left[1-\frac{2}{3}\beta\!\left(\frac{r_1}{r_2}\right)^{\!1/2}\right] \notag\\[4pt]
&\simeq \frac{3}{2}\,\dot{M}_0\!\left(\frac{M}{r_1}-\frac{M}{r_2}\right)
   \quad\text{for}\;r_2 \gg r_1 \gg r_1.
\tag{5.2.20}
\end{align}

It is instructive to examine the origin of this heat energy by asking at what
net rate energy is being deposited in the region $r_1 < r < r_2$. Gravitational
potential energy is released at a rate $\dot{M}_0(M/r_1 - M/r_2)$, but only half of this
energy can go into orbital kinetic energy
(virial theorem). Thus, the rate at which gravitational energy gets converted
into heat is
\[
\tfrac{1}{2}\dot{M}_0(M/r_1 - M/r_2).
\]

The viscous stresses transport outward not only angular momentum, but also
energy. The rate at which energy is transported across radius $r$ is
\[
\dot{E} = \dot{\varOmega}J
= \dot{\varOmega}(2\pi r\cdot 2h\cdot t_\phi^{\;r})\bigl[1-\Omega(r/r')^{1/2}\bigr].
\]
Thus, energy is deposited by viscosity between radii $r_1$ and $r_2$ at a rate
\[
\dot{M}_0\!\left[\frac{M}{r_1}\left(1-\beta\!\left(\frac{r_1}{r_2}\right)^{\!1/2}\right)
-\frac{M}{r_2}\left(1-\beta\!\left(\frac{r_1}{r_2}\right)^{\!1/2}\right)\right].
\]

This energy deposition rate plus the deposition rate for gravitational energy is
equal to the total heating rate. Notice that at radii $r \gg r_1$, gravity accounts
directly for only one-third of the heat. The remaining two-thirds is transported
into the heating region by viscous stresses. Much of the early literature on disk
accretion, e.g., Lynden-Bell (1969)\cite{LyndenBell1969}, failed to take account of energy transport
by viscous stresses, and therefore underestimated by a factor 3 the heating at
$r \gg r_1$.

If none of the heat were radiated away, the thermal energy of the disk
would be 3/2 times its gravitational potential energy; and its temperature
would thus be
\[
T \sim (M/r) m_p \sim (10^{13}\,\text{K})(r/10^5\,\text{cm})^{-1}(M/M_\odot).
\]

Such temperatures are absurdly high. Thermal bremsstrahlung and other
radiative processes will never permit such temperatures to arise. Instead, they
will quickly remove almost all of the heat that is generated. This means that

%% ============================================================
%% Page 417, right column
%% ============================================================

the total flux $2F$ from the top and bottom faces of the disk must equal the
heating rate per unit area, as calculated above (the heat flows out of the disk
vertically rather than radially because the disk is so thin):
\begin{equation}
F = \frac{3\dot{M}_0 M}{8\pi r^2\,r_1}
    \left[1-\beta\!\left(\frac{r_1}{r}\right)^{\!1/2}\right].
\tag{5.2.19}
\end{equation}

The total power radiated is thus
\begin{equation}
L = \int_{r_1}^{\infty} 2F\cdot 2\pi r\,dr
  = \bigl(\tfrac{3}{2}-\beta\bigr)\dot{M}_0\,M/r_1\,.
\tag{5.2.20}
\end{equation}
Of this total power radiated, $\frac{1}{2}\dot{M}_0 M/r_1$ comes from gravity, while
$(1-\beta)\dot{M}_0 M/r_1$ comes from the rotational energy of the star.

\paragraph*{Spectrum}

A detailed treatment of the spectrum radiated will be given in \S5.10. Here we
only note that, if the radiation is blackbody, then the surface temperature of
the disk at radius $r$ must be
\begin{equation}
T_s = \left(\frac{4F}{b}\right)^{\!1/4}
\simeq (3\times 10^7\,\text{K})
\left(\frac{\dot{M}_0}{10^{-9}\,M_\odot/\text{yr}}\right)^{\!1/4}
\!\left(\frac{M}{M_\odot}\right)^{\!-1/2}
\!\left(\frac{M}{r}\right)^{\!3/4}
\left[1-\beta\!\left(\frac{r_1}{r}\right)^{\!1/2}\right]^{1/4}.
\tag{5.2.21}
\end{equation}

Because most of the radiation is emitted from $r/M \sim 10$ for a black hole or
neutron star, and from $r/M \sim 10^5$ for a dwarf, and because the black-body
spectrum peaks at a photon energy of
\begin{equation}
h\nu_{\max} \simeq (2.44\times 10^{-4}\;\text{eV})(T_s/{}^\circ\!\text{K}),
\tag{5.2.22}
\end{equation}
the spectrum of the total radiation from the disk should peak at
\begin{alignat}{2}
h\nu_{\max} &\simeq (1\;\text{keV})
\left(\frac{\dot{M}_0}{10^{-9}\,M_\odot/\text{yr}}\right)^{\!1/4}
\!\left(\frac{M}{M_\odot}\right)^{\!-1/2}
&\quad&\text{for neutron star or hole} \notag\\[4pt]
h\nu_{\max} &\simeq (0.01\;\text{keV})
\left(\frac{\dot{M}_0}{10^{-9}\,M_\odot/\text{yr}}\right)^{\!1/4}
&\quad&\text{for white dwarf.}
\tag{5.2.23}
\end{alignat}

Thus, for reasonable accretion rates the disk can emit strongly in the X-ray
region (${\sim}\,1$ to 10~keV); but the X-ray spectrum should fall fairly rapidly with
increasing energy. It actually turns out that electron-scattering opacity impedes
the emission of blackbody radiation, and causes the spectrum to peak at energies
higher than that (5.2.23). (See \S5.10.) However, the above estimates are still
accurate to within a factor ${\sim}\,10$.

%% ============================================================
%% Page 418, left column: Vertical pressure balance
%% ============================================================

\paragraph*{Vertical pressure balance}

The thickness of the disk is governed by a balance between vertical pressure force
and the tidal gravitational force of the compact star (``vertical momentum
conservation''):
\begin{equation}
\frac{dp}{dz} = \rho_0 \times
\bigl(\text{``acceleration of gravity''}\bigr) = \rho_0\,\frac{Mz}{r^3}\,.
\tag{5.2.24}
\end{equation}
The approximate solution to this equation is
\begin{equation}
h \simeq (p/\rho_0)^{1/2}(r^3/M)^{1/2} = c_s/\Omega.
\tag{5.2.25}
\end{equation}
Here $h$ is the half-thickness of the disk, $c_s$ is the speed of sound in the gas, and
$\Omega = (M/r^3)^{1/2}$ is the angular velocity of the gas.

\paragraph*{Sources of viscosity}

All of the above conclusions are based on conservation laws alone. To proceed
further, one must make some assumption about the magnitude of the viscosity.
The dominant sources of viscosity are probably chaotic magnetic fields, and
turbulence in the gas flow.

For cases of interest the normal star (typically a B0 supergiant) could
have a surface magnetic field of $B_s \sim 100$~gauss. Field lines will be dragged
by the flowing matter off of the normal star and into the disk. The deposited field
should be rather chaotic because there is no preferred direction for the field
in the gas that flows off the normal star. Turbulence, if present in the disk,
will also make the field chaotic. Once the chaotic field has been deposited in
the disk, the shear of the flow will magnify it at a rate (eq.~(2.5.20))
\begin{equation}
dB_\phi/dr = \alpha_\phi B_r,\qquad dB_\phi/dt = 0
\tag{5.2.26}
\end{equation}
corresponding to an increase of $B_\phi$ by the amount $B_r$ with every circuit around
the compact star. This growth of field will be counterbalanced, at least in part,
by reconnection of field lines at the interfaces between chaotic cells (see \S2.9),
and perhaps also by a bulging of field lines out of the disk, pinch-off of field
lines, and escape of ``magnetic bubbles.'' From a macroscopic viewpoint (scale
large compared to chaotic cells) the effects of the chaotic field can be described by
pressure and viscosity. The magnetic pressure of the field will be of order $B^2/8\pi$, where $B$
is the mean field strength. Because the shearing of the chaotic field will tend to
string it out along the $\phi$-direction, its shear stress $t_{\phi r}$ will be somewhat (perhaps
a factor~$10^2$) smaller than its pressure
\begin{equation}
t_{\phi r}^{(\text{mag})} \simeq \nu_0\,\rho_0\,r\,\frac{d\Omega}{dr}
= B^2/8\pi.
\tag{5.2.27}
\end{equation}
The magnetic pressure cannot exceed the thermal pressure
\begin{equation}
p^{(\text{mag})} \simeq p^{(\text{therm})} \simeq \rho_0 c_s^2;
\tag{5.2.28}
\end{equation}

%% ============================================================
%% Page 419, left column
%% ============================================================

otherwise the field lines would bulge out of the disk, reconnect, and escape
(``bubbles''). Thus, the magnetic viscous stresses will satisfy
\begin{equation}
t_{\phi r}^{(\text{mag})} \lesssim p = \rho_0 c_s^2.
\tag{5.2.29}
\end{equation}
(Here $c_s$ is the speed of sound, and $p$ is the total pressure.)

The gas flow in the disk may well be turbulent. The coefficient of dynamic
viscosity associated with turbulence is
\begin{equation}
\eta \sim \rho_0\, v_{\text{turb}}\, l_{\text{turb}},
\tag{5.2.30}
\end{equation}
where $v_{\text{turb}}$ is the speed of the turbulent motions relative to the mean rest
frame of the gas, and $l_{\text{turb}}$ is the characteristic size of the largest turbulent cells
[see, e.g., \S31 of Landau and Lifshitz~(1959)\cite{LandauLifshitz1959}].
If the turbulent speed ever exceeds
the sound speed, then shocks develop and quickly convert the turbulent energy
into heat. Thus, $v_{\text{turb}} \lesssim c_s$. The turbulent scale is limited by the disk thickness,
$l_{\text{turb}} \lesssim h$. Consequently, the shear stress due to turbulence is bounded by
\begin{equation}
t_{\phi r}^{(\text{turb})} \sim \eta\,\rho_0\,r\,\frac{d\Omega}{dr}
\lesssim (\nu_0\,\rho_0\,c_s\,h)\,\Omega
\simeq \rho_0 c_s^2 \simeq p.
\tag{5.2.31}
\end{equation}
Inequalities (5.2.29) and (5.2.31) on the magnetic and turbulent stresses are
identical. They suggest that one dump one's lack of knowledge about the true
magnitude of the viscosity into a single parameter $\alpha$ defined by
\begin{equation}
t_{\phi r} = \alpha p.
\tag{5.2.32}
\end{equation}
Someday, perhaps ten years hence, when one understands the magnetic and
turbulent viscosities better, one can insert into the formalism a reliable value of
$\alpha$. In the meantime one only knows that
\begin{equation}
\alpha \lesssim 1.
\tag{5.2.33}
\end{equation}
If $\alpha \sim 1$, the disk will have a rather mottled structure on scales of order~$h$; if
$\alpha \ll 1$ it will be rather smooth on such scales.

\paragraph*{Radiative transport}

The heat generated by viscosity must be transported vertically to the surface of
the disk before it can be radiated. The disk turns out to be optically thick
($\tau_{ff}+\tau_{es}\gg 1$ at $z=0$). Hence, one calculates the energy transport using the
diffusion approximation [eq.~(2.6.43) reduced to Newtonian form]:
\begin{equation}
\frac{d}{dz}\!\left(b\,T^4\right) = \bar{\kappa}\,\rho_0\,\frac{d^2}{dz^2}\,F.
\tag{5.2.34}
\end{equation}
The approximate solution to this equation of transport is
\begin{equation}
bT^4 \simeq \bar{\kappa}\;\Sigma\; F.
\tag{5.2.35}
\end{equation}
